% =============================================================================
% CAPÍTULO 1 — INTRODUCCIÓN Y DISPOSICIONES GENERALES
% =============================================================================
\chapter[Introducción y Generalidades]{Introducción y Disposiciones Generales}%
\label{ch:introduccion}
\resumenCapitulo{Este capítulo establece el propósito, el alcance y la audiencia del
Manual de Instalación de la plataforma EthosHub. Identifica con precisión las versiones
de los componentes de software a desplegar, describe la naturaleza del artefacto
monolítico integrado y fija el vocabulario controlado y la normativa que rigen el resto
del documento. Constituye el marco de referencia que el operador debe leer antes de
iniciar cualquier procedimiento técnico.}

Este Manual de Instalación describe, paso a paso, el procedimiento completo para
construir, configurar, desplegar y verificar la plataforma \textbf{EthosHub} en un
entorno de servidor. Se redacta conforme al estándar \textbf{ISO/IEC/IEEE 15289:2024}
—que define el contenido de los productos de información del ciclo de vida del software—
y, en su tratamiento de requisitos, conforme a \textbf{IEEE 830}. A diferencia del
Manual Técnico, que documenta el \textit{diseño} interno del sistema, este documento es
de naturaleza \textbf{operativa}: su objetivo es que un operador con conocimientos
medios pueda llevar el producto desde el código fuente hasta un servicio en
funcionamiento, sin ambigüedad y de forma reproducible.\par

\begin{AvisoNota}
Lea este capítulo en su totalidad antes de ejecutar cualquier comando. Las versiones
exactas de las herramientas (capítulo~\ref{ch:requisitos}) y la correcta provisión de
variables de entorno (capítulo~\ref{ch:preparacion}) condicionan el éxito de todo el
proceso. Saltarse la verificación de pre-requisitos es la causa más frecuente de fallos
de despliegue descritos en el capítulo~\ref{ch:troubleshooting}.
\end{AvisoNota}

% -----------------------------------------------------------------------------
\section{Propósito del Documento}%
\label{sec:proposito}

El propósito de este manual es servir como \textbf{guía autoritativa y reproducible}
del despliegue de EthosHub en su versión \texttt{2.0.0}. En términos del estándar
ISO/IEC/IEEE 15289:2024, constituye un \textit{procedimiento} (\textit{procedure}):
una secuencia ordenada de pasos verificables que conducen a un resultado determinado.
El documento persigue cuatro objetivos concretos:\par

\begin{enumerate}[noitemsep]
    \item \textbf{Reproducibilidad.} Garantizar que dos operadores distintos, partiendo
    del mismo código fuente y siguiendo estos pasos, obtengan exactamente el mismo
    artefacto desplegado y el mismo comportamiento del servicio.
    \item \textbf{Trazabilidad.} Vincular cada instrucción a un artefacto real del
    repositorio (\comando{build.sh}, \comando{pom.xml}, \comando{Dockerfile},
    \comando{.env.production}), de modo que toda afirmación sea verificable.
    \item \textbf{Autonomía del operador.} Permitir que el personal de infraestructura
    complete la instalación sin necesidad de consultar al equipo de desarrollo,
    anticipando los errores comunes y sus mitigaciones.
    \item \textbf{Cumplimiento normativo.} Documentar el despliegue de forma auditable,
    con control de versiones, identificación de roles y matriz de trazabilidad, según
    exige la convocatoria CPTIS-2302-2026.
\end{enumerate}

Quedan \textbf{fuera del alcance} de este documento el diseño detallado del software
(cubierto por el Manual Técnico), la operativa funcional para el usuario final
(cubierta por el Manual de Usuario) y la administración continua del proyecto Supabase
en la nube, que se asume gestionada por el equipo de base de datos. El presente manual
abarca desde la obtención del código fuente hasta la verificación de un servicio
operativo, incluyendo su desinstalación y reversión controlada.\par

% -----------------------------------------------------------------------------
\section{Alcance del Sistema}%
\label{sec:alcance}

\subsection{Descripción General de EthosHub (Talent Hub)}%
\label{subsec:descripcion}

\textbf{EthosHub} es una plataforma unificada de \textit{Talent Hub} que articula dos
experiencias complementarias sobre un mismo dominio de datos. Por un lado, el
\textbf{profesional} construye un portafolio público enriquecido —proyectos con
documentos PDF e importación desde Google Drive, experiencia laboral, formación
académica, habilidades y un currículum interactivo generado en el \textit{CV Studio}
con asistencia de inteligencia artificial—. Por otro, el \textbf{reclutador} descubre,
filtra y conecta con talento mediante un \textit{pipeline} tipo Kanban, métricas, un
\textit{leaderboard} e \textit{insights} asistidos por IA.\par

Desde la perspectiva de la instalación, lo relevante es que ambas experiencias se
sirven desde un \textbf{único artefacto ejecutable}. El sistema no se compone de
múltiples servicios que haya que orquestar por separado: el backend Spring Boot empaqueta
en su interior la aplicación React ya compilada y la entrega como contenido estático.
La tabla~\ref{tab:stack-instalacion} resume el \textit{stack} tecnológico que el operador
encontrará durante el despliegue, organizado por capa.\par

\vspace{0.2cm}
\noindent
\renewcommand{\arraystretch}{1.35}
\fontsize{8.7}{10.4}\selectfont\sffamily
\begin{tabularx}{\dimexpr\textwidth-2pt\relax}{|>{\bfseries\color{colorPrimario}}l|l|X|}
\hline
\rowcolor{headerTabla}
\textbf{Capa} & \textbf{Tecnología} & \textbf{Relevancia para la instalación} \\
\hline
Lenguaje (BE) & Java 17 & Requiere JDK 17+ instalado y en el \texttt{PATH}. \\
\hline
\rowcolor{grisTecnico}
Framework (BE) & Spring Boot 4.0.6 & Provee Actuator (\texttt{/actuator/health}) y el JAR ejecutable. \\
\hline
Persistencia & jOOQ + Spring JDBC & Acceso SQL tipado; exige consistencia esquema/clases generadas. \\
\hline
\rowcolor{grisTecnico}
Base de datos & PostgreSQL 15+ (Supabase) & Servicio externo; se aprovisiona, no se instala localmente. \\
\hline
Lenguaje (FE) & TypeScript 5.3 (\textit{strict}) & Se compila con \texttt{tsc} durante el \textit{build}. \\
\hline
\rowcolor{grisTecnico}
Framework (FE) & React 18.3 + Vite 5.1 & Genera el directorio \texttt{dist/} que se inyecta en el JAR. \\
\hline
IA & Google Gemini (REST) & Servicio externo; requiere \texttt{GEMINI\_API\_KEY}. \\
\hline
\rowcolor{grisTecnico}
Documentación API & springdoc-openapi 2.7 & Expone Swagger UI para la verificación del capítulo~\ref{ch:verificacion}. \\
\hline
Empaquetado & Maven + Docker & Generan el JAR y la imagen de contenedor. \\
\hline
\end{tabularx}\normalfont\normalsize%
\label{tab:stack-instalacion}

\par\vspace{0.3cm}
\begin{NotaConfiguracion}
EthosHub \textbf{no aloja su propia base de datos ni su propio servicio de
autenticación}: ambos se delegan en \textbf{Supabase} (PostgreSQL gestionado, Auth y
Realtime). En consecuencia, la instalación no incluye levantar un PostgreSQL local en
producción, sino \textbf{conectar} el artefacto a un proyecto Supabase ya provisto. El
único PostgreSQL local que aparece en este manual es el de \textbf{Testcontainers},
empleado exclusivamente durante las pruebas automatizadas (apartado~\ref{sec:testcontainers}).
\end{NotaConfiguracion}

\subsection{Arquitectura del Artefacto Monolítico Integrado}%
\label{subsec:arquitectura}

El patrón de despliegue de EthosHub es un \textbf{monolito empaquetado}
(\textit{packaged monolith}): la SPA de React se compila a archivos estáticos
(\comando{dist/}) que se copian al directorio \ruta{src/main/resources/static/} del
backend y se sirven desde el mismo JAR ejecutable que expone la API REST. El resultado
es un \textbf{único artefacto} (\comando{ethoshub-0.0.1-SNAPSHOT.jar}) que, al
arrancar, atiende tanto las peticiones de la interfaz como las de la API en el mismo
origen (\textit{same-origin}). La figura~\ref{fig:monolito} ilustra esta composición.\par

\vspace{0.3cm}
\begin{figure}[h!]
\centering
\begin{tikzpicture}[node distance=0.7cm and 1.1cm]
    % Límite del JAR
    \node[draw=c4BordeLim, dashed, thick, rounded corners=8pt,
          minimum width=12.5cm, minimum height=6.2cm, label={[font=\sffamily\footnotesize\bfseries\color{c4Sistema}]above:JAR monolítico \texttt{ethoshub.jar} (puerto 8080)}] (jar) at (0,0) {};

    \node[c4componente, fill=c4Componente!50] (spa) at (-3.6,1.4) {SPA React (\texttt{dist/})\\\scriptsize servida en \texttt{/}};
    \node[c4componente] (fallback) at (-3.6,-0.1) {\texttt{SpaFallbackFilter}\\\scriptsize entrega \texttt{index.html}};
    \node[c4componente] (sec) at (-3.6,-1.5) {Spring Security\\\scriptsize Resource Server (JWT)};

    \node[c4contenedor] (ctrl) at (1.2,1.4) {Controllers REST\\\scriptsize \texttt{/api}, \texttt{/v1}, \texttt{/auth}};
    \node[c4contenedor, fill=c4Contenedor!85] (svc) at (1.2,-0.1) {Services};
    \node[c4contenedor, fill=c4Contenedor!70] (repo) at (1.2,-1.5) {Repositories\\\scriptsize jOOQ / \texttt{DSLContext}};

    \draw[c4flecha] (ctrl) -- (svc);
    \draw[c4flecha] (svc) -- (repo);
    \draw[c4flecha] (sec) -- (ctrl);
    \draw[c4flecha] (fallback) -- (spa);

    % Externos
    \node[c4datos, right=1.4cm of repo] (db) {PostgreSQL\\\scriptsize Supabase};
    \node[c4externo, above=0.5cm of db] (gemini) {Google\\Gemini};
    \draw[c4flecha] (repo) -- (db);
    \draw[c4flecha] (svc.east) to[bend left=12] (gemini.west);
\end{tikzpicture}
\caption{Composición del artefacto monolítico: la SPA compilada y la API REST conviven en un único JAR que se conecta a servicios externos (Supabase, Gemini).}%
\label{fig:monolito}
\end{figure}

Esta arquitectura tiene tres consecuencias prácticas de gran importancia para el
operador, que conviene interiorizar desde el inicio:\par

\begin{itemize}[noitemsep]
    \item \textbf{Un solo proceso que desplegar.} En producción solo se ejecuta el JAR
    (o su contenedor Docker). No hay servidor web separado para el \textit{frontend} ni
    necesidad de configurar CORS entre dominios, pues todo es \textit{same-origin}.
    \item \textbf{El \textit{frontend} debe compilarse \textit{antes} del empaquetado.}
    Si el directorio \comando{dist/} no se genera correctamente, el JAR se empaqueta sin
    interfaz. El script \comando{build.sh} (capítulo~\ref{ch:construccion}) automatiza y
    ordena estas fases.
    \item \textbf{En desarrollo, la topología es distinta.} Durante la ingeniería local,
    el \textit{frontend} corre en su propio servidor Vite (puerto~3000) y reenvía las
    llamadas de API al backend (puerto~8080) mediante un \textit{proxy}. Esta diferencia
    se detalla en el apartado~\ref{sec:dev-mode}.
\end{itemize}

% -----------------------------------------------------------------------------
\section{Identificación del Producto y Versión}%
\label{sec:identificacion}

Conforme a ISO/IEC/IEEE 15289:2024, todo procedimiento de instalación debe identificar
sin ambigüedad el producto y la versión que despliega. EthosHub adopta
\textbf{Versionado Semántico} (\textit{SemVer}, \texttt{MAYOR.MENOR.PARCHE}). La versión
documentada en esta guía es la \texttt{2.0.0}, integrada por dos componentes de
repositorio independientes que se ensamblan en el artefacto final.\par

\subsection{Componente Backend (\texttt{ethoshubB} v2.0.0)}%
\label{subsec:backend-id}

El repositorio \comando{ethoshubB} contiene el backend Spring Boot y es el responsable
de generar el artefacto final. Sus señas de identidad para la instalación se resumen en
la tabla~\ref{tab:id-backend}.\par

\vspace{0.2cm}
\noindent
\renewcommand{\arraystretch}{1.35}
\fontsize{9}{10.8}\selectfont\sffamily
\begin{tabularx}{\dimexpr\textwidth-2pt\relax}{|>{\bfseries\color{colorPrimario}}l|X|}
\hline
\rowcolor{headerTabla}
\multicolumn{2}{|l|}{\sffamily\textbf{Componente Backend}} \\
\hline
Identificador del repositorio & \texttt{ethoshubB} \\
\hline
\rowcolor{grisTecnico}
Coordenadas Maven & \texttt{com.ethoshub:ethoshub:0.0.1-SNAPSHOT} \\
\hline
Ruta local de referencia & \ruta{/home/dpardo/Documentos/ethoshubB} \\
\hline
\rowcolor{grisTecnico}
Artefacto generado & \texttt{target/ethoshub-0.0.1-SNAPSHOT.jar} \\
\hline
Punto de entrada & \texttt{com.ethoshub.EthosHubApplication} \\
\hline
\rowcolor{grisTecnico}
Puerto por defecto & \texttt{8080} \\
\hline
\end{tabularx}\normalfont\normalsize%
\label{tab:id-backend}

\subsection{Componente Frontend (\texttt{ethoshubF} v2.0.0)}%
\label{subsec:frontend-id}

El repositorio \comando{ethoshubF} contiene la SPA en React + TypeScript. No genera un
artefacto desplegable por sí mismo: produce el directorio \comando{dist/} que el backend
absorbe. Sus señas de identidad se resumen en la tabla~\ref{tab:id-frontend}.\par

\vspace{0.2cm}
\noindent
\renewcommand{\arraystretch}{1.35}
\fontsize{9}{10.8}\selectfont\sffamily
\begin{tabularx}{\dimexpr\textwidth-2pt\relax}{|>{\bfseries\color{colorPrimario}}l|X|}
\hline
\rowcolor{headerTabla}
\multicolumn{2}{|l|}{\sffamily\textbf{Componente Frontend}} \\
\hline
Identificador del repositorio & \texttt{ethoshubF} \\
\hline
\rowcolor{grisTecnico}
Nombre del paquete (npm) & \texttt{ethoshub@1.0.0} \\
\hline
Ruta local de referencia & \ruta{/home/dpardo/Documentos/ethoshubF} \\
\hline
\rowcolor{grisTecnico}
Artefacto generado & \texttt{dist/} (estáticos) $\rightarrow$ \ruta{ethoshubB/src/main/resources/static/} \\
\hline
Servidor de desarrollo & Vite, puerto \texttt{3000} \\
\hline
\rowcolor{grisTecnico}
Idiomas (i18next) & Español (\texttt{es}, fuente de verdad), Inglés (\texttt{en}), Portugués (\texttt{pt}) \\
\hline
\end{tabularx}\normalfont\normalsize%
\label{tab:id-frontend}

\par\vspace{0.3cm}
\begin{AvisoPrecaucion}
La relación de rutas entre ambos repositorios es \textbf{rígida}: el script de
construcción asume que \comando{ethoshubF} es un directorio hermano de \comando{ethoshubB}
(\comando{../ethoshubF}). Si los clona en una estructura distinta, deberá ajustar la
variable \comando{frontend.dir} del \comando{pom.xml} y la ruta \comando{FRONTEND\_DIR}
de \comando{build.sh}, según se detalla en el apartado~\ref{sec:clonado}.
\end{AvisoPrecaucion}

% -----------------------------------------------------------------------------
\section{Audiencia Objetivo y Niveles de Competencia Técnica}%
\label{sec:audiencia}

Este manual está orientado, principalmente, a la \textbf{persona que realiza la
instalación}: un operador de infraestructura, un evaluador técnico de la convocatoria o
un desarrollador que prepara un entorno de pruebas. No se asume experiencia previa con
EthosHub, pero sí un nivel básico de competencia en línea de comandos. La
tabla~\ref{tab:audiencia} detalla los perfiles, el nivel mínimo esperado y las secciones
prioritarias para cada uno.\par

\vspace{0.2cm}
\noindent
\renewcommand{\arraystretch}{1.4}
\fontsize{8.7}{10.4}\selectfont\sffamily
\begin{tabularx}{\dimexpr\textwidth-2pt\relax}{|>{\bfseries\color{colorPrimario}}l|X|c|}
\hline
\rowcolor{headerTabla}
\textbf{Perfil} & \textbf{Competencias mínimas asumidas} & \textbf{Caps.} \\
\hline
Operador de despliegue & Terminal Linux, Docker básico, edición de archivos \texttt{.env}. & 2--5, 9 \\
\hline
\rowcolor{grisTecnico}
Evaluador técnico (TIS) & Lectura de logs, navegador web, conceptos de cliente/servidor. & 1, 5, 7 \\
\hline
Administrador de BD & Supabase, PostgreSQL, SQL, políticas RLS. & 3, 6 \\
\hline
\rowcolor{grisTecnico}
Usuario evaluador (PWA) & Uso de navegador, ajustes de privacidad, cámara/micrófono. & 6.3 \\
\hline
\end{tabularx}\normalfont\normalsize%
\label{tab:audiencia}

\par\vspace{0.3cm}
Para mantener el documento accesible, cada procedimiento se presenta como una
\textbf{secuencia numerada de pasos}, acompañada de los comandos literales que deben
ejecutarse y del resultado esperado. Los conceptos que exceden el nivel del operador
(por ejemplo, el funcionamiento interno de las políticas RLS) se delimitan en cajas de
\textit{Nota de Configuración} y pueden delegarse al administrador de base de datos.\par

% -----------------------------------------------------------------------------
\section{Glosario de Términos y Acrónimos}%
\label{sec:glosario}

El glosario de la tabla~\ref{tab:glosario} fija el vocabulario controlado que se emplea
de forma recurrente en los procedimientos. Se recomienda al operador familiarizarse con
estos términos antes de continuar, pues se utilizan sin redefinición en los capítulos
siguientes.\par

\vspace{0.2cm}
\noindent
\renewcommand{\arraystretch}{1.4}
\fontsize{9}{10.8}\selectfont\sffamily
\begin{tabularx}{\dimexpr\textwidth-2pt\relax}{|>{\bfseries\color{colorPrimario}}l|X|}
\hline
\rowcolor{headerTabla}
\textbf{Término} & \textbf{Definición operativa} \\
\hline
PWA & \textit{Progressive Web App}. Aplicación web que ofrece capacidades cercanas a las nativas (instalación, acceso a cámara/micrófono). En EthosHub habilita la captura de vídeo/audio del \textit{CV Studio}; exige origen seguro (HTTPS o \texttt{localhost}), restricción central del capítulo~\ref{ch:seguridad}. \\
\hline
\rowcolor{grisTecnico}
SPA & \textit{Single Page Application}. Cliente React que carga una sola página y reescribe su contenido por JavaScript. Es lo que se compila a \texttt{dist/} y se sirve desde el JAR. \\
\hline
JWT & \textit{JSON Web Token}. Credencial firmada que Supabase emite tras la autenticación y que el backend valida en cada petición protegida. \\
\hline
\rowcolor{grisTecnico}
JWKS & \textit{JSON Web Key Set}. Conjunto de claves públicas publicado por Supabase (\texttt{/auth/v1/.well-known/jwks.json}) con el que el backend verifica la firma de los JWT. \\
\hline
RLS & \textit{Row Level Security}. Mecanismo de PostgreSQL que restringe el acceso a filas según el usuario. Se configura en Supabase (apartado~\ref{sec:rls}). \\
\hline
\rowcolor{grisTecnico}
jOOQ & Biblioteca de acceso a datos que genera clases Java tipadas a partir del esquema SQL. Una desalineación esquema/clases produce el error del apartado~\ref{sec:error-jooq}. \\
\hline
Monolito empaquetado & Patrón de despliegue en el que la SPA y la API se sirven desde un único JAR. Ver figura~\ref{fig:monolito}. \\
\hline
\rowcolor{grisTecnico}
Resource Server & Componente OAuth2 que protege recursos validando JWT sin emitirlos. Rol del backend de EthosHub. \\
\hline
Artefacto & Producto de la construcción listo para desplegar: en este manual, el JAR ejecutable o la imagen Docker. \\
\hline
\end{tabularx}\normalfont\normalsize%
\label{tab:glosario}

% -----------------------------------------------------------------------------
\section{Documentos de Referencia y Normativa Aplicable}%
\label{sec:referencias}

La elaboración de esta guía se rige por los estándares y documentos de la
tabla~\ref{tab:referencias}. El operador no necesita haberlos leído para completar la
instalación, pero su cumplimiento explica la estructura y el nivel de detalle del
documento.\par

\vspace{0.2cm}
\noindent
\renewcommand{\arraystretch}{1.35}
\fontsize{8.7}{10.4}\selectfont\sffamily
\begin{tabularx}{\dimexpr\textwidth-2pt\relax}{|>{\bfseries\color{colorPrimario}}l|X|}
\hline
\rowcolor{headerTabla}
\textbf{Documento} & \textbf{Aplicación en este manual} \\
\hline
ISO/IEC/IEEE 15289:2024 & Define el contenido del procedimiento de instalación: identificación, control de versiones, roles y trazabilidad. \\
\hline
\rowcolor{grisTecnico}
IEEE 830-1998 & Marco para la expresión de requisitos (capítulo~\ref{ch:requisitos}): hardware, software y red. \\
\hline
Manual Técnico de EthosHub (\texttt{MT-ARQ-ETHOSHUB-2.0.0}) & Documento hermano que describe el diseño interno; se referencia para detalles arquitectónicos. \\
\hline
\rowcolor{grisTecnico}
\texttt{README} de \texttt{ethoshubB} y \texttt{ethoshubF} & Fuente primaria de los comandos y variables reproducidos en los capítulos~\ref{ch:construccion} y~\ref{ch:preparacion}. \\
\hline
Convocatoria CPTIS-2302-2026 & Marco contractual del proyecto y origen de los requisitos de documentación. \\
\hline
\end{tabularx}\normalfont\normalsize%
\label{tab:referencias}

\par\vspace{0.3cm}
Con el marco de referencia establecido, el capítulo~\ref{ch:requisitos} detalla los
requisitos previos de hardware, software y red que deben satisfacerse antes de obtener el
código fuente.\par
